% Budget: ~1 page. Measured part: logs.txt section 17 (src/diag_qiea_profile.py,
% results/diag_qiea_profile*.txt). Framing follows submission2.txt section 8
% ("profiling-guided performance engineering") -- confirm with co-authors.
\section{Runtime Profile and Parallel Design}
\label{sec:parallel}

Before parallelizing anything, we measured where a run spends its time. On a
laptop with an Intel Core i7-12700H and BLAS pinned to one thread, we timed full 500-generation
runs (3 seeds, medians) on a small, a medium and a large instance, and split the
wall-clock time into three parts: fitness evaluation (route split and the five
objectives), archive maintenance, and everything else, i.e., guide selection,
variation operators and neighborhood replacement. We used low-overhead timers for
the shares and a separate cProfile run only to rank functions, since cProfile
inflates pure-Python code.

\begin{table}[t]
\centering
\caption{Runtime breakdown of one HQIEA-ARGC run (500 generations, median of 3
seeds). $S_{\max}$: Amdahl bound on speedup if only evaluation is parallelized.}
\label{tab:profile}
\setlength{\tabcolsep}{3.5pt}
\begin{tabular}{@{}lrrrrrr@{}}
\toprule
Instance & $n$ & Time & Eval. & Archive & Other & $S_{\max}$ \\
\midrule
A-n32-k5    & 31  & 7.7\,s  & 29\% & 2\% & 68\% & 1.4$\times$ \\
route2\_199 & 198 & 16.2\,s & 63\% & 2\% & 35\% & 2.7$\times$ \\
route1\_334 & 333 & 24.2\,s & 73\% & 2\% & 25\% & 3.7$\times$ \\
\bottomrule
\end{tabular}
\end{table}

\Cref{tab:profile} shows the result. Evaluation takes 29--73\% of the runtime and
its share grows with instance size, but each call is cheap: 0.03--0.27\,ms, so a
whole generation's batch of 126 evaluations costs only 4--34\,ms. Archive
maintenance is negligible at about 2\%. The remaining time is almost constant
across instances (5.2--6.0\,s per run) and comes mostly from one source: the
Tchebycheff function is called about 1.9 million times per run (500 generations
$\times$ 126 subproblems $\times$ 10 neighbors $\times$ 3 calls), each time on a
5-element array. The cost is interpreter and library call overhead, not
arithmetic. The restart adds 0--5\% to the runtime, all of it from the extra
evaluations it triggers.

These numbers bound what parallel evaluation alone can achieve. By Amdahl's law,
even with unlimited workers the speedup cannot exceed 1.4$\times$ on the small
instance and 3.7$\times$ on the largest (\cref{tab:profile}). Moreover, the
per-generation batch is small enough that distributing it across processes risks
costing more in communication than it saves. The profile therefore suggests a
different order of work:
\begin{enumerate}
\item \emph{Remove the serial overhead.} Guide selection can compute all
  Tchebycheff values with a single array operation per generation. Replacement is
  sequential, because a neighbor's solution can change within the loop, but the
  child's scores can be computed in advance and the current scores kept in an
  array, which preserves the exact semantics. Evaluation can be compiled with
  Numba.
\item \emph{Parallelize at the right level.} Within one run, evaluation of a batch
  is independent across children, and because all random numbers are drawn in the
  main process, a parallel run with the same seed should produce the same result
  as a serial one; we verify this before timing anything. We compare process-based (multiprocessing) and thread-based (Numba
  \texttt{prange}) parallel evaluation. Across runs, the 1{,}050 independent runs
  of the full study (7 instances $\times$ 5 algorithms $\times$ 30 seeds) are
  embarrassingly parallel.
\item \emph{Account for the hybrid CPU.} The i7-12700H has 6 performance cores
  (12 threads) and 8 efficiency cores. Speedup curves are expected to bend once
  work spills onto the efficiency cores, so we measure performance-core-only,
  efficiency-core-only and mixed placements.
\end{enumerate}

% TODO: results of steps 1-3 go in Sec. experiments (runtime subsection).
% Do NOT state speedups here until measured.
